% ─── 中文摘要 ───
\begin{chineseabstract}
无人机搭载电磁超声探头进行结构无损检测时，探头与被测表面的接触状态难以精确判定。传统方法依赖单一超声回波信号的幅值阈值判断，在飞行扰动和姿态变化条件下易产生误判。针对这一核心问题，本文提出了一种基于多模态逐级物理约束的无人机自主电磁超声探测框架，构建了融合电磁超声回波序列、无人机位姿动态残差及RGBD局部视觉ROI特征的时序感知模型。

在理论层面，本文首先从电磁-弹性耦合机理出发，将EMAT信号视为由激励、传播、反射和接收四个物理环节串联形成的可观测动态过程，推导了洛伦兹体力作为弹性动力学方程体载荷项的基本关系。在此基础上，对接触边界和升距效应进行了回波统计表征建模，将接触状态定义为由升距、法向约束、末端力学稳定性和回波可观测性共同决定的多因素隐变量，并以概率形式表达其连续演化。

在平台与数据层面，本文完成了无人机多模态实验平台的搭建与集成。系统以多旋翼无人机为载体，搭载Livox MID-360激光雷达、Intel RealSense D435 RGBD相机、PX4飞控与MAVROS通信链路，采用FAST-LIO 2.0提供实时六自由度位姿估计。开发了基于CH346C芯片的EMAT探头ROS驱动节点，实现了约40 Hz的稳定波形采集与Qt5可视化。构建了多模态同步数据采集系统，可同步记录位姿、RGBD视频流与EMAT回波信号，为模型训练提供结构化数据支撑。

在算法层面，针对EMAT回波（约40 Hz）、飞控里程计（约100 Hz）与RGBD视觉（约30 Hz）的异步采样问题，提出了基于统一时间基准的加权插值对齐与独立编码器映射方案。在此基础上，设计了物理约束注意力机制，将时间邻近约束和空间残差约束以对数偏置形式嵌入Transformer的缩放点积注意力计算中，使跨模态交互同时服从时间连续性、空间邻近性和运动学合理性。进一步引入了模态可靠性门控权重和双阈值持续时间约束状态机，实现了从传感器不确定性评估到控制闭环滞回判别的完整证据链。

仿真与实验结果表明，本文提出的多模态物理约束融合框架能够有效克服单模态阈值方法的局限性，在临界接触段实现状态判别概率的平滑过渡，为无人机自主电磁超声检测提供了从物理建模、信号采集、特征融合到闭环控制的一体化解决方案。

\chinesekeyword{无人机无损检测，电磁超声，多模态融合，物理约束注意力，接触状态判别，Transformer}
\end{chineseabstract}
